\pageheading{5 An optional causal statement for actual receive tiles}
\textbf{Move the spatial exclusion into one reusable theorem only if its conclusion has value to later clients.} The current \code{far_west_zero} and \code{far_east_zero} proofs independently inspect the neighboring wall, plain tile or head transition. A stronger statement packages the reason behind both arguments: every receive tile inside a live encoded row is located exactly at the actual moved head.

Let $p=\mathrm{cfg.head}$, let $a$ be the actual transition at $p$, and let $n=a.\mathrm{move.apply}(p)$. The proposed \code{receive_at_next_head} assumes allowed tiles, horizontal matching, \code{PersistentStrip}, the original \code{RowEncodes}, a live state, and an actual directional receive tuple at an interior position $x$. It concludes $x=n$. It has \textbf{no temporal-matching or no-finish hypothesis}. Its target assumption is membership in \code{targets d}, which retains target 215.

\textbf{The argument is spatial.}
\begin{enumerate}
\item A receive tile has a nonzero incoming rail, at least 6. Inspect its incoming-side neighbor, one integer cell away.
\item That neighbor is either interior or exactly the wall. A wall is impossible because its facing edge is at most 2 on the left or 4 on the right.
\item If it is an interior non-head cell, both north inputs are plain. The existing \code{plain_pair_edge_zero} forces the matching rail to zero, a contradiction.
\item The neighbor is therefore the actual head. Its forced transition has one inactive zero edge. Matching the positive rail identifies the move direction, and the adjacent-coordinate equation gives $x=n$.
\end{enumerate}
This proof needs positivity and the inactive edge, not rail injectivity. Exact target identification remains owned by the existing directional receiver theorems; it is not claimed as an additional result of this helper.

The second helper, \code{far_pass_or_finish}, eliminates only the receive case of the exact plain-input grammar. Away from $p$ and $n$, with $b$ the original tape bit, it concludes
\[
 T(x,t)=\langle0,b,0,b\rangle\quad\lor\quad
 T(x,t)=\langle0,b,0,438\rangle.
\]
Both zero-edge clients become projections of this disjunction. Only \code{far_pass} invokes the independently derived \code{no_south438} result to remove the finish alternative. This keeps the temporal premise exactly where it belongs.

\textbf{Cost and decision.} The two edge proofs shrink from 22 and 21 proof lines to two each. But the new causal and disjunction proofs cost 49 and 7 lines. Including the unchanged four-line \code{far_pass}, the affected bodies grow from \textbf{47 to 64 lines}; the full file grows from \textbf{205 to 252 physical lines}. Thus this is an optional semantic API, not a current net-shortening recommendation. Prefer it only when future spatial-only consumers need receiver location, or when centralizing this invariant is worth the source cost.

The seven existing exported statements remain unchanged. The supplied candidate is concrete Lean source, not pseudocode that assumes a future direction API. Its sources are \src{WangTiles942RowStep}{35}{122}, \src{WangTiles942Signals}{54}{104}, and the full new declarations in \code{row/helper.lean}. The novelty is this exact row-level causal conclusion; constructor views and direction abstraction are already credited to the baseline.
